\par\Needspace{6\baselineskip}\medskip
\hypertarget{native-112-article083}{}\pdfbookmark[2]{83.}{article-083}
\begin{center}83.\end{center}
Will man bei der Lösung derselben Aufgabe mit der Berechnung des Winkels $\Pi$ beginnen, so verfährt man wie folgt. Wir gehen von den Gleichungen I aus und ziehen die dritte von der zweiten, die dritte von der ersten und die zweite von der ersten ab. So erhalten wir die folgenden drei neuen Gleichungen.
\[
\begin{aligned}
&\frac{\dfrac1{r'}-\dfrac1{r''}}{2\sin\tfrac12(N''-N')}=\frac{e}{p}\sin(\tfrac12N'+\mathpunct{\cdot}\tfrac12N''-\Pi)\\[6pt]
\mathrm{II.}\qquad&\frac{\dfrac1r-\dfrac1{r''}}{2\sin\tfrac12(N''-N)}=\frac{e}{p}\sin(\tfrac12N+\tfrac12N''-\Pi)\\[6pt]
&\frac{\dfrac1r-\dfrac1{r'}}{2\sin\tfrac12(N'-N)}=\frac{e}{p}\sin(\tfrac12N+\tfrac12N'-\Pi).
\end{aligned}
\]
Je zwei dieser Gleichungen liefern nach Lemma II aus Art.\ 78 die Größen $\Pi$ und $\dfrac{e}{p}$. Daraus erhält man mit einer beliebigen der Gleichungen I auch $e$ und $p$. Wendet man die dritte in Art.\ 78, II angegebene Lösung an, so erhält man durch Kombination der ersten mit der dritten der Gleichungen II den folgenden Algorithmus. Man bestimme den Hilfswinkel $\zeta$ aus der Gleichung
\[
\operatorname{tang}\zeta=\frac{\dfrac{r'}r-1}{1-\dfrac{r'}{r''}}\cdot\frac{\sin\tfrac12(N''-N')}{\sin\tfrac12(N'-N)}\mathpunct{\raisebox{.45ex}{,}}
\]
\SourceLeaf{113}dann gilt
\[
\operatorname{tang}(\tfrac14N+\tfrac12N'+\tfrac14N''-\Pi)=\operatorname{tang}(45^0+\zeta)\operatorname{tang}\tfrac14(N''-N).
\]
Vertauscht man den zweiten Ort mit dem ersten oder dem dritten, so erhält man zwei weitere, ganz analoge Lösungen. Da dieses Verfahren weniger bequeme Formeln für $\dfrac{e}{p}$ liefert, bestimmt man $e$ und $p$ besser aus zwei der Gleichungen I nach dem Verfahren aus Art.\ 80. Die Mehrdeutigkeit bei der Bestimmung von $\Pi$ aus dem Tangens des Winkels $\tfrac14N+\tfrac12N'+\tfrac14N''-\Pi$ ist so aufzulösen, dass $e$ positiv wird. Offenbar ergeben sich für $e$ entgegengesetzte Werte, wenn man für $\Pi$ Werte wählt, die sich um $180^0$ unterscheiden. Das Vorzeichen von $p$ ist hingegen von dieser Mehrdeutigkeit unabhängig. Der Wert von $p$ kann nur dann negativ werden, wenn die drei gegebenen Punkte auf dem von der Sonne abgewandten Teil der Hyperbel liegen. Dieser Fall widerspricht den Naturgesetzen und bleibt hier außer Betracht.

Was sich bei Anwendung des ersten Verfahrens aus Art.\ 78, II nach aufwendigeren Substitutionen ergeben würde, lässt sich im vorliegenden Fall bequemer auf folgende Weise gewinnen. Man multipliziere die erste der Gleichungen II mit $\cos\tfrac12(N''-N')$, die dritte mit $\cos\tfrac12(N'-N)$ und ziehe das zweite Produkt vom ersten ab. Bei entsprechender Anwendung von Lemma I aus Art.\ 78\footnote{Dazu setzt man in der zweiten Formel $A=\tfrac12(N''-N')$, $B=\tfrac12N+\tfrac12N''-\Pi$, $C=\tfrac12(N-N'')$.} ergibt sich dann die Gleichung
\[
\begin{aligned}
&\tfrac12\left(\dfrac1{r'}-\dfrac1{r''}\right)\operatorname{cotang}\tfrac12(N''-N')-\tfrac12\left(\dfrac1r-\dfrac1{r'}\right)\operatorname{cotang}\tfrac12(N'-N)\\[4pt]
&\qquad=\frac{e}{p}\sin\tfrac12(N''-N)\cos(\tfrac12N+\tfrac12N''-\Pi).
\end{aligned}
\]
Verbindet man diese Gleichung mit der zweiten der Gleichungen II, so erhält man $\Pi$ und $\dfrac{e}{p}$, und zwar $\Pi$ nach der Formel
\[
\begin{aligned}
&\operatorname{tang}(\tfrac12N+\tfrac12N''-\Pi)\\[4pt]
&\qquad=\frac{\dfrac{r'}r-\dfrac{r'}{r''}}{\left(1-\dfrac{r'}{r''}\right)\operatorname{cotang}\tfrac12(N''-N')-\left(\dfrac{r'}r-1\right)\operatorname{cotang}\tfrac12(N'-N)}\mathpunct{\cdot}
\end{aligned}
\]
Auch daraus lassen sich zwei weitere, ganz analoge Formeln ableiten, indem man den zweiten Ort mit dem ersten oder dem dritten vertauscht.
